% Part: first-order-logic
% Chapter: models-of-arithmetic
% Section: models-of-pa

\documentclass[../../../include/open-logic-section]{subfiles}

\begin{document}

\olfileid{mod}{mar}{mpa}
\section{$\Th{PA}$-র মডেল}

\begin{explain}
$\Th{TA}$-র যেকোনো অমানক মডেল $\Th{PA}$-রও মডেল। আমরা জানি,
$\Th{TA}$-র, এবং তাই $\Th{PA}$-রও, অমানক মডেল আছে। আরও জানি, এমন
অমানক মডেলে অমানক ``সংখ্যা'' থাকে; অর্থাৎ সংজ্ঞাক্ষেত্রের এমন
!!{element}s থাকে, যেগুলি সব মানক ``সংখ্যা''-র ``পরে''। কিন্তু তারা
কীভাবে বিন্যস্ত? তাদের সংখ্যা কত? আমরা দেখেছি, দুর্বলতর তত্ত্ব
$\Th{Q}$-এর মডেলে একটি মাত্র অমানক সংখ্যাও থাকতে পারে। কিন্তু ওই সরল
!!{structure}s-গুলি $\Th{PA}$ বা $\Th{TA}$-র মডেল নয়।

$\Th{PA}$ বা $\Th{TA}$-র মডেলের গঠন বোঝার চাবিকাঠি হলো, এই
তত্ত্বগুলিতে কোন তথ্যগুলি !!{derivable} তা দেখা। যেমন, $\Th{PA}$
ইতিমধ্যেই $\lforall[x][\eq/[x][x']]$ এবং
$\lforall[x][\lforall[y][\eq[(x+y)][(y+x)]]]$ প্রমাণ করে; ফলে যেসব
সরল গঠনে এই !!{sentence}s মিথ্যা, সেগুলি $\Th{PA}$-র মডেল হতে পারে না।

ধরা যাক, $\Struct{M}$ হলো $\Th{PA}$-র একটি মডেল। তাহলে
$\Th{PA} \Proves !A$ হলে $\Sat{M}{!A}$। আবার
$\Assign{\Obj{0}}{M}$-এর জন্য $\nszero$, $\Assign{\prime}{M}$-এর জন্য
$\nssucc$, $\Assign{+}{M}$-এর জন্য $\nsplus$,
$\Assign{\times}{M}$-এর জন্য $\nstimes$, এবং $\Assign{<}{M}$-এর জন্য
$\nsless$ ব্যবহার করি। তখন যেকোনো !!{sentence}~$!A$, $\nszero$,
$\nssucc$, $\nsplus$, $\nstimes$ ও~$\nsless$ সম্পর্কে কোনো শর্ত বলে;
আর $\Sat{M}{!A}$ হলে সেই শর্তটি পরিতৃপ্ত হতে হবে। যেমন,
$\Sat{M}{!Q_1}$, অর্থাৎ
$\Sat{M}{\lforall[x][\lforall[y][(\eq[x'][y'] \lif \eq[x][y])]]}$ হলে
$\nssucc$-কে !!{injective} হতে হবে।
\end{explain}

\begin{prop}
$\Struct{M}$-এ $\nsless$ একটি কঠোর রৈখিক ক্রম; অর্থাৎ এটি নিচের
শর্তগুলি পরিতৃপ্ত করে:
\begin{enumerate}
\item কোনো~$x \in \Domain{M}$-এর জন্যই $x \nsless x$ নয়।
\item $x \nsless y$ এবং $y \nsless z$ হলে $x \nsless z$।
\item যেকোনো $x \neq y$-এর জন্য $x \nsless y$ অথবা $y \nsless x$।
\end{enumerate}
\end{prop}

\begin{proof}
$\Th{PA}$ নিচেরগুলি প্রমাণ করে:
\begin{enumerate}
\item $\lforall[x][\lnot x < x]$
\item $\lforall[x][\lforall[y][\lforall[z][((x < y \land y < z) \lif x < z)]]]$
\item $\lforall[x][\lforall[y][((x < y \lor y < x) \lor \eq[x][y])]]$
\end{enumerate}
% BN-SRC-128 TARGET-CORRECTION-BEGIN
\par\smallskip\noindent\textnormal{\small\textbf{উৎস-সংশোধন (BN-SRC-128):} হিমায়িত ইংরেজি উৎসে তৃতীয় সর্বব্যাপী সূত্রটির শেষে খোলা গোলবন্ধনীর তুলনায় একটি অতিরিক্ত সমাপ্তি-গোলবন্ধনী আছে। বাংলা সূত্রে ওই অতিরিক্ত চিহ্নটি বাদ দিয়ে ত্রিবিভাগের সূত্রটি ভারসাম্যপূর্ণ করা হয়েছে।} % BN-SRC-128 TARGET-CORRECTION-END
\end{proof}

\begin{prop}
\ollabel{prop:M-discrete} $\nsless$-ক্রমে $\nszero$ হলো
$\Domain{M}$-এর ক্ষুদ্রতম !!{element}। যেকোনো~$x$-এর জন্য
$x \nsless x^\nssucc$, এবং এই ধর্মসম্পন্ন $\nsless$-ক্ষুদ্রতম
!!{element} হলো $x^\nssucc$। শূন্য নয় এমন যেকোনো~$x$-এর জন্য এমন
একটি অনন্য~$y$ আছে,
যাতে $y^\nssucc=x$। (আমরা $y$-কে $\Struct{M}$-এ $x$-এর
``পূর্বসূরি'' বলি এবং $^\nssucc x$ দিয়ে চিহ্নিত করি।)
\end{prop}
% BN-SRC-129 TARGET-CORRECTION-BEGIN
\par\smallskip\noindent\textnormal{\small\textbf{উৎস-সংশোধন (BN-SRC-129):} হিমায়িত প্রতিজ্ঞাটি প্রত্যেক~$x$-এর একটি পূর্বসূরি দাবি করে, অথচ একই প্রতিজ্ঞায় $\nszero$-কে ক্ষুদ্রতম উপাদান বলা হয়েছে এবং $\Th{PA}$-তে শূন্য কোনো উত্তরসূরি নয়। বাংলা বিবৃতি ও পরের খণ্ড-প্রমাণে পূর্বসূরির দাবিটি শূন্য নয় এমন উপাদানে সীমিত করা হয়েছে; অমানক খণ্ডের সব উপাদান এই শর্ত পূরণ করে।} % BN-SRC-129 TARGET-CORRECTION-END

\begin{proof}
অনুশীলনী।
\end{proof}

\begin{prob}
$\Th{PA}$-তে !!{derivable} (এবং তাই $\Struct{N}$-এ সত্য) এমন
$\Lang{L_A}$-!!{sentence}s খুঁজে বার করো, যেগুলি
\olref[mod][mar][mpa]{prop:M-discrete}-এ $\nszero$, $\nssucc$ ও
$\nsless$-এর ধর্মগুলি নিশ্চিত করে।
\end{prob}

\begin{prop}
$\Struct{M}$-এর সব মানক !!{element}s, সব অমানক !!{element}s-এর চেয়ে
($\nsless$ অনুসারে) ছোট।
\end{prop}

\begin{proof}
$\Struct{M}$-এর মানক !!{element}~$\Value{\num{n}}{M}$-এর সংক্ষেপে
$n$ লিখব। যেকোনো~$n \in \Nat$-এর জন্য $\Th{Q}$ ইতিমধ্যেই প্রমাণ করে
$\lforall[x][(x < \num{n}' \lif (\eq[x][\num{0}] \lor
  \eq[x][\num{1}] \lor \dots \lor \eq[x][\num{n}]))]$। কোনো
এমন কোনো !!{element}s নেই, যা $\nsless \nszero$। তাই $n$ মানক ও~$x$ অমানক হলে
$x \nsless n$ হতে পারে না। সংজ্ঞা অনুসারে অমানক উপাদান এমন উপাদান,
যা কোনো~$n \in \Nat$-এর জন্য $\Value{\num{n}}{M}$ নয়; অতএব
$x \neq n$-ও। $\nsless$ রৈখিক ক্রম বলে তখন $n \nsless x$ হতেই হবে।
\end{proof}

\begin{prop}
$\Domain{M}$-এর প্রত্যেক অমানক !!{element}~$x$ নিচের উপসেটটির উপাদান:
\[
\dots ^{\nssucc\nssucc\nssucc}x \nsless ^{\nssucc\nssucc}x \nsless
^{\nssucc}x \nsless x \nsless x^{\nssucc} \nsless x^{\nssucc\nssucc}
\nsless x^{\nssucc\nssucc\nssucc} \nsless \dots
\]
এই উপসেটকে $x$-এর \emph{খণ্ড} বলি এবং $[x]$ লিখি। এর কোনো ক্ষুদ্রতম
বা বৃহত্তম !!{element} নেই। একে সেইসব $y \in \Domain{M}$-এর সেট
হিসেবে চিহ্নিত করা যায়, যাদের জন্য কোনো মানক~$n$-এ
$x \nsplus n=y$ অথবা $y \nsplus n=x$।
\end{prop}

\begin{proof}
এমন সেট~$[x]$ সব সময়ই আছে; কারণ $\Domain{M}$-এর প্রত্যেক
!!{element}~$y$-এর একটি অনন্য উত্তরসূরি~$y^\nssucc$ আছে এবং শূন্য নয়
এমন প্রত্যেকটির একটি অনন্য পূর্বসূরি~$^\nssucc y$ আছে। পরপর
!!{element}s~$y$, $y^\nssucc$-এর জন্য $y \nsless y^\nssucc$, এবং
$y^\nssucc$ হলো $\nsless$-ক্রমে $\Domain{M}$-এর ক্ষুদ্রতম
!!{element}, যার চেয়ে $y$ হলো $\nsless$-ছোট। সব সময়
$^\nssucc y \nsless y$ এবং
$y \nsless y^\nssucc$ বলে $[x]$-এর কোনো ক্ষুদ্রতম বা বৃহত্তম
!!{element} নেই। $y \in [x]$ হলে $x \in [y]$; কারণ তখন হয়
$y^{\nssucc\dots\nssucc}=x$, নয় $x^{\nssucc\dots\nssucc}=y$। যদি
$y^{\nssucc\dots\nssucc}=x$ হয় ($n$-টি $\nssucc$-সহ), তবে
$y \nsplus n=x$, এবং বিপরীতটিও সত্য; কারণ ($n$ যদি $\prime$-এর সংখ্যা
হয়) $\Th{PA} \Proves \lforall[x][\eq[x^{\prime\dots\prime}][(x +
    \num{n})]]$।
\end{proof}

\begin{prop}
$[x] \neq [y]$ এবং $x \nsless y$ হলে, প্রত্যেক $u \in [x]$ ও প্রত্যেক
$v \in [y]$-এর জন্য $u \nsless v$।
\end{prop}

\begin{proof}
লক্ষ করো, $\Th{PA} \Proves \lforall[x][\lforall[y][(x < y \lif (x' < y
    \lor x' = y))]]$। অতএব $u \nsless v$ হলে এবং $[u]\neq[v]$ হলে,
যেকোনো~$n$-এর জন্য $u \nsplus n^\nssucc \nsless v$-ও হয়।

যেকোনো $u \in [x]$ হলো $\nsless y$: অনুমান অনুসারে
$x \nsless y$। $u \nsless x$ হলে সংক্রমণশীলতায় $u \nsless y$। আর
$x \nsless u$ কিন্তু $u \in [x]$ হলে কোনো~$n$-এর জন্য
$u=x\nsplus n^\nssucc$; তাই এইমাত্র প্রমাণিত তথ্য থেকে $u \nsless y$।

এবার ধরা যাক, $v \in [y]$ হলো $\nsless y$; অর্থাৎ কোনো মানক~$m$-এর
জন্য $v \nsplus m^\nssucc=y$। এতে $v \nsless x$ অসম্ভব; তা হলে
$y=v \nsplus m^\nssucc \nsless x$ হতো। আবার স্পষ্টতই $x\neq v$;
নইলে $x \nsplus m^\nssucc=v \nsplus m^\nssucc=y$ হয়ে
$[x]=[y]$ হতো। অতএব $x \nsless v$। কিন্তু তখন যেকোনো~$n$-এর জন্য
$x \nsplus n^\nssucc \nsless v$-ও। ফলে $x \nsless u$ এবং
$u \in [x]$ হলে $u \nsless v$। $u \nsless x$ হলে সংক্রমণশীলতায়ও
$u \nsless v$।

সব শেষে, $y \nsless v$ হলে $u \nsless v$; কারণ আমরা দেখিয়েছি
$u \nsless y$, আর $y \nsless v$।
\end{proof}

\begin{cor}
$[x] \neq [y]$ হলে $[x] \cap [y] = \emptyset$।
\end{cor}

\begin{proof}
ধরা যাক, $z \in [x]$ এবং $x \nsless y$। তখন সব $u \in [y]$-এর জন্য
$z \nsless u$। যদি $z \in [y]$ হতো, তবে $z \nsless z$ পেতাম।
$y \nsless x$ হলেও একই যুক্তি।
\end{proof}

\begin{explain}
এর অর্থ, খণ্ডগুলিকেই এমনভাবে ক্রমিত করা যায় যা $\nsless$-কে মানে:
$[x] \nsless [y]$ যদি এবং কেবল যদি $x \nsless y$; অথবা সমতুল্যভাবে,
যেকোনো $u \in [x]$ ও~$v \in [y]$-এর জন্য $u \nsless v$ হলে। স্পষ্টতই
মানক খণ্ড~$[0]$ ক্ষুদ্রতম খণ্ড। কোনো অমানক খণ্ডের সঙ্গে তার ছেদ নেই,
এবং দুটি অমানক খণ্ডের মধ্যেও ছেদ নেই। বিশেষত বারবার উত্তরসূরি বা
পূর্বসূরি নিয়ে অন্য কোনো খণ্ডে ``পৌঁছানো'' যায় না।
\end{explain}

\begin{prop}
$x$ ও~$y$ অমানক হলে $x \nsless x \nsplus y$ এবং
$x \nsplus y \notin [x]$।
\end{prop}

\begin{proof}
$y$ অমানক হলে $y \neq \nszero$। $\Th{PA} \Proves
\lforall[x][(y \neq \Obj{0} \lif x < (x+y))]$। এবার ধরা যাক,
$x \nsplus y \in [x]$। যেহেতু $x \nsless x \nsplus y$, তাই
$x \nsplus n^\nssucc=x \nsplus y$ হতো। কিন্তু
$\Th{PA} \Proves
\lforall[x][\lforall[y][\lforall[z][(\eq[(x+y)][(x+z)] \lif y = z)]]]$
(যোগের কর্তন-বিধি)। এর অর্থ হতো, কোনো মানক~$n$-এর জন্য
$y=n^\nssucc$; কিন্তু $y$-কে অমানক ধরে নেওয়া হয়েছে।
\end{proof}

\begin{prop}
কোনো ক্ষুদ্রতম অমানক খণ্ড নেই।
\end{prop}

\begin{proof}
$\Th{PA} \Proves \lforall[x][\lexists[y][(\eq[(y+y)][x] \lor
      \eq[(y+y)'][x])]]$; অর্থাৎ প্রত্যেক $x$-কে $2$ দিয়ে ভাগ করা যায়
(সম্ভবত ভাগশেষ~$1$ রেখে)। $x$ অমানক হলে $y$-ও অমানক। আগের প্রতিজ্ঞা
অনুসারে $y \nsless y \nsplus y$ এবং $y \nsplus y \notin [y]$। তখন
$y \nsless (y \nsplus y)^\nssucc$ এবং
$(y \nsplus y)^\nssucc \notin [y]$-ও। কিন্তু হয়
$x=y \nsplus y$, নয় $x=(y \nsplus y)^\nssucc$; তাই $y \nsless x$
এবং $y \notin [x]$।
\end{proof}

\begin{prop}
কোনো বৃহত্তম খণ্ড নেই।
\end{prop}

\begin{proof}
অনুশীলনী।
\end{proof}

\begin{prob}
দেখাও যে $\Th{PA}$-র কোনো অমানক মডেলে বৃহত্তম খণ্ড নেই।
\end{prob}

\begin{prop}
\ollabel{prop:blocks-dense}
খণ্ডগুলির ক্রম ঘন। অর্থাৎ $x \nsless y$ এবং $[x] \neq [y]$ হলে
দুটির থেকে ভিন্ন এমন একটি খণ্ড~$[z]$ আছে, যা তাদের মাঝখানে।
\end{prop}

\begin{proof}
ধরা যাক, $x \nsless y$। আগের মতো $x \nsplus y$-কে দুই দিয়ে ভাগ করা
যায় (সম্ভবত ভাগশেষ রেখে): এমন $z \in \Domain{M}$ আছে, যাতে হয়
$x \nsplus y = z \nsplus z$, নয়
$x \nsplus y = (z \nsplus z)^\nssucc$। উপাদান~$z$, $x$ ও~$y$-এর
``গড়'', এবং $x \nsless z$ ও~$z \nsless y$।
% BN-SRC-130 TARGET-CORRECTION-BEGIN
\par\smallskip\noindent\textnormal{\small\textbf{উৎস-সংশোধন (BN-SRC-130):} এই অংশে গঠনের যোগের জন্য সর্বত্র $\nsplus$ স্থির করা হলেও হিমায়িত উৎসের গড়-সমীকরণে চারবার অনির্ধারিত $\oplus$ ব্যবহৃত। বাংলা সূত্রে ওই চারটি চিহ্নকে $\nsplus$ করা হয়েছে; সমীকরণের অন্য সব অংশ অপরিবর্তিত।} % BN-SRC-130 TARGET-CORRECTION-END
\end{proof}

\begin{prob}
\olref[mod][mar][mpa]{prop:blocks-dense}-এর বিস্তারিত প্রমাণ লেখো।
$z$-এর অস্তিত্ব নিশ্চিত করতে $\Th{PA}$-কে কোন !!{sentence}
!!{derive} করতে হবে? কেন $x \nsless z$ ও~$z \nsless y$, এবং কেন
$[x] \neq [z]$ ও~$[z] \neq [y]$?
\end{prob}

\begin{explain}
অতএব কোনো গণনীয় অমানক মডেলের অমানক খণ্ডগুলি মূলদ সংখ্যার মতো
ক্রমিত: তারা অন্তবিন্দুহীন একটি !!a{denumerable} ঘন রৈখিক ক্রম গঠন করে।
% BN-SRC-131 TARGET-CORRECTION-BEGIN
\par\smallskip\noindent\textnormal{\small\textbf{উৎস-সংশোধন (BN-SRC-131):} হিমায়িত উৎসটি কোনো গণনীয়তার অনুমান ছাড়াই সব অমানক খণ্ডকে $\Rat$-এর মতো এবং !!{denumerable} বলে, যদিও অগণনীয় মডেলের খণ্ডসমষ্টিও অগণনীয় হতে পারে। বাংলা পাঠে এই সিদ্ধান্তকে গণনীয় অমানক মডেলে সীমিত করা হয়েছে; ঘনতা ও অন্তবিন্দুহীনতার আগের দাবিগুলি অপরিবর্তিত।} % BN-SRC-131 TARGET-CORRECTION-END
দেখানো যায়, এমন যেকোনো দুটি !!{denumerable} ক্রম পরস্পর সমরূপ। এর
ফলে সত্য পাটিগণিতের যেকোনো দুটি !!{enumerable} অমানক মডেল
$\Struct{M}_1$ ও~$\Struct{M_2}$-এর কেবল $<$ ও~$=$-যুক্ত ভাষায়
রিডাক্টগুলি সমরূপ। বস্তুত সমরূপতা~$h$ নিচের মতো সংজ্ঞায়িত করা যায়:
$\Struct{M_1}$ ও~$\Struct{M_2}$-এর মানক অংশগুলি মানক মডেল
$\Struct{N}$-এর সমরূপ, তাই পরস্পরেরও সমরূপ। অমানক অংশ গঠনকারী
খণ্ডগুলি নিজেরাই মূলদ সংখ্যার মতো ক্রমিত, তাই সেগুলিও সমরূপ। প্রতিটি
খণ্ডে নির্বিচারে উপাদান মিলিয়ে খণ্ডগুলির সমরূপতাকে খণ্ডের
\emph{ভিতরে} বিস্তৃত করা যায়; তারপর $\Struct{M_1}$-এ $x$-এর
উত্তরসূরির প্রতিচ্ছবি নিই $\Struct{M_2}$-এ $x$-এর প্রতিচ্ছবির
উত্তরসূরি। লক্ষ করো, এর থেকে পূর্ণ পাটিগণিতের ভাষায়
$\mathfrak{M}_1$ ও~$\mathfrak{M}_2$ সমরূপ---এ কথা \emph{অনুসৃত হয় না}
(সমরূপতা সব সময়ই কোনো !!a{language}-এর সাপেক্ষে); কারণ
$\Domain{M_1}$ ও~$\Domain{M_2}$-এর উপর যোগ ও গুণ সংজ্ঞায়িত করার
অসমরূপ উপায় আছে। (ভটের একটি বিখ্যাত উপপাদ্য থেকেও এটি অনুসৃত হয়:
কোনো সম্পূর্ণ তত্ত্বের গণনীয় মডেলের সংখ্যা~2 হতে পারে না।)
\end{explain}
\end{document}
